\section{Experiments}
\label{sec:experiments}

The experiments evaluate \method{} through four result views.  The feature
prediction experiment is the main comparison across Web relation, program,
circuit, timing, molecular, and social graphs.  The autoregressive experiment
uses the same point set object for rule guided graph expansion.  The ablation
experiment measures how content recovery, relation recovery, structural
deviation, training perturbation, and set reading contribute to the final
performance.  The WL collision experiment studies whether recovery trained point
sets separate controlled nonisomorphic graph pairs from BREC and WL diagnostic
sets \cite{morris2019wl,wang2023brec}.  Together, these four experiments test
task utility, graph growth, component attribution, and controlled structural
separation.

\subsection{Feature Prediction Main Experiment}
\label{sec:feature-main}

The main experiment evaluates supervised prediction on heterogeneous graph
tasks.  Amazon Photo, Coauthor CS, and Cora represent product co-purchase,
academic collaboration, and citation networks
\cite{shchur2018,sen2008,yang2016}.  HiGraph evaluates hierarchical program
control flow graphs for malware classification \cite{chen2025higraph}.  OCB and
TimingPredict evaluate analog circuit performance and pre-routing timing
prediction \cite{dong2023cktgnn,guo2022timing}.  AIDS and IMDB evaluate
molecular and social graph classification from standard graph learning
benchmarks \cite{morris2020,yang2015}.  These tasks cover node labels, graph
labels, and continuous graph properties, so the learned point set is tested
under several downstream readout forms.

The baselines cover attribute only prediction, neighborhood aggregation,
propagation, long range graph sequence modeling, and graph to point set
modeling.  MLP uses attributes without relation input.  GCN, GraphSAGE, GAT, and
GIN represent widely used neighborhood computation mechanisms
\cite{kipf2017,hamilton2017,velickovic2018,xu2019}.  APPNP adds personalized
PageRank propagation after prediction \cite{gasteiger2019}.  Graph Mamba models
long range graph dependencies through selective state spaces
\cite{wang2024graphmamba}.  PST maps graphs into point set representations
\cite{wang2024graphpointset}.  Table~\ref{tab:feature-results} reports the
embedding score and the final task result as separate cells for each method.
Across the ten listed comparisons, \method{} gives the best final task result in
every row, showing that the recoverable point set can serve as an effective
supervised representation.

\begin{table*}[t]
  \centering
  \caption{Feature prediction main results from Experiment Result Table A. Each method keeps Emb and Result as separate columns.}
  \label{tab:feature-results}
  \tiny
  \setlength{\tabcolsep}{1.1pt}
  \resizebox{\textwidth}{!}{%
  \begin{tabular}{llcccccccccccccccccccccc}
    \toprule
    Dataset & Metric
    & \multicolumn{2}{c}{MLP}
    & \multicolumn{2}{c}{GCN}
    & \multicolumn{2}{c}{GraphSAGE}
    & \multicolumn{2}{c}{GAT}
    & \multicolumn{2}{c}{GIN}
    & \multicolumn{2}{c}{APPNP}
    & \multicolumn{2}{c}{SGFormer}
    & \multicolumn{2}{c}{Polynormer}
    & \multicolumn{2}{c}{Graph-Mamba}
    & \multicolumn{2}{c}{PST}
    & \multicolumn{2}{c}{\method{}} \\
    \cmidrule(lr){3-4}\cmidrule(lr){5-6}\cmidrule(lr){7-8}
    \cmidrule(lr){9-10}\cmidrule(lr){11-12}\cmidrule(lr){13-14}
    \cmidrule(lr){15-16}\cmidrule(lr){17-18}\cmidrule(lr){19-20}
    \cmidrule(lr){21-22}\cmidrule(lr){23-24}
    & & Emb & Result & Emb & Result & Emb & Result & Emb & Result
    & Emb & Result & Emb & Result & Emb & Result & Emb & Result
    & Emb & Result & Emb & Result & Emb & Result \\
    \midrule
    HiGraph & F1 & 0.52 & 0.45 & 0.67 & 0.62 & 0.52 & 0.60 & 0.52 & 0.59 & 0.61 & 0.51 & 0.56 & 0.40 & 0.57 & 0.41 & 0.57 & 0.41 & 0.56 & 0.40 & 0.61 & 0.55 & \textbf{0.81} & \textbf{0.86} \\
    OCB & r & 0.68 & 0.70 & 0.71 & 0.69 & 0.71 & 0.69 & 0.72 & 0.71 & 0.71 & 0.70 & 0.68 & 0.62 & 0.67 & 0.60 & 0.68 & 0.61 & 0.65 & 0.56 & 0.70 & 0.68 & \textbf{0.84} & \textbf{0.90} \\
    Timing & R\textsuperscript{2} & 0.74 & 0.52 & 0.85 & 0.68 & 0.83 & 0.66 & 0.79 & 0.64 & 0.83 & 0.65 & 0.86 & 0.70 & 0.74 & 0.74 & 0.75 & 0.76 & 0.74 & 0.73 & 0.77 & 0.71 & \textbf{1.00} & \textbf{0.95} \\
    Photo & Acc & 0.80 & 0.85 & 0.87 & 0.91 & 0.79 & 0.91 & 0.84 & 0.88 & 0.78 & 0.87 & 0.85 & 0.90 & 0.88 & 0.94 & 0.86 & 0.95 & 0.88 & 0.94 & 0.85 & 0.92 & \textbf{0.88} & \textbf{0.95} \\
    Coauthor & Acc & 0.86 & 0.94 & 0.87 & 0.94 & 0.88 & 0.94 & 0.87 & 0.94 & 0.86 & 0.94 & 0.89 & 0.93 & 0.87 & 0.95 & 0.82 & 0.94 & 0.86 & 0.94 & 0.87 & 0.94 & \textbf{0.90} & \textbf{0.95} \\
    Photo int. & Acc & 0.76 & 0.80 & 0.82 & 0.88 & 0.80 & 0.90 & 0.82 & 0.84 & 0.78 & 0.83 & 0.86 & 0.88 & 0.88 & 0.93 & 0.86 & 0.94 & 0.88 & 0.93 & 0.85 & 0.91 & \textbf{0.97} & \textbf{0.94} \\
    Cora & Acc & 0.43 & 0.47 & 0.84 & 0.75 & 0.73 & 0.70 & 0.78 & 0.73 & 0.70 & 0.65 & 0.87 & 0.80 & 0.81 & 0.66 & 0.80 & 0.73 & 0.86 & 0.69 & 0.81 & 0.76 & \textbf{0.89} & \textbf{0.82} \\
    AIDS & Acc & 0.67 & 0.68 & 0.77 & 0.80 & 0.76 & 0.80 & 0.76 & 0.80 & 0.76 & 0.80 & 0.75 & 0.75 & 0.76 & 0.77 & 0.76 & 0.78 & 0.75 & 0.76 & 0.78 & 0.84 & \textbf{0.89} & \textbf{0.99} \\
    IMDB-M & Acc & 0.51 & 0.35 & 0.54 & 0.36 & 0.52 & 0.35 & 0.52 & 0.34 & 0.56 & 0.40 & 0.53 & 0.35 & 0.54 & 0.35 & 0.52 & 0.35 & 0.50 & 0.35 & 0.55 & 0.41 & \textbf{0.63} & \textbf{0.52} \\
    IMDB-B & Acc & 0.55 & 0.43 & 0.61 & 0.50 & 0.61 & 0.50 & 0.61 & 0.50 & 0.61 & 0.51 & 0.60 & 0.47 & 0.61 & 0.48 & 0.61 & 0.49 & 0.60 & 0.48 & 0.62 & 0.55 & \textbf{0.75} & \textbf{0.75} \\
    \bottomrule
  \end{tabular}}
\end{table*}

\subsection{Autoregressive Graph Modeling}
\label{sec:ar-results}

The autoregressive experiment evaluates graph expansion with the same set
representation used for prediction.  During rollout, the set expansion module
adds live points, the relation decoder turns the updated point set into a graph,
and the legality signal selects states that satisfy the target rule.  The three
rules test complementary structural requirements: planar validity captures a
global embeddability constraint, chordal validity captures the required cycle
structure, and bipartite validity captures a two part relation constraint.

The comparison includes recurrent, attention based, flow based, score based, and
diffusion based graph generation methods, including GraphRNN, GRAN, GraphAF,
GDSS, and DiGress \cite{you2018,liao2019,shi2020,jo2022,vignac2023}.  The full
workbook records both embedding scores and terminal validity.  To keep the
paper table compact, Table~\ref{tab:autoregressive-results} reports terminal
validity, the key rollout metric.  \method{} obtains the strongest validity on
all four rollout settings, which supports the use of point set expansion as a
graph growth interface.

\begin{table}[t]
  \centering
  \caption{Autoregressive validity from Experiment Result Table B.}
  \label{tab:autoregressive-results}
  \scriptsize
  \setlength{\tabcolsep}{3pt}
  \begin{tabular}{lcc}
    \toprule
    Task & Best baseline & \method{} \\
    \midrule
    Planar Legal & 0.85 & \textbf{0.94} \\
    Planar Net & 0.85 & \textbf{0.97} \\
    Chordal & 0.24 & \textbf{0.83} \\
    Bipartite & 0.18 & \textbf{0.74} \\
    \bottomrule
  \end{tabular}
\end{table}

\subsection{Ablation Experiment}
\label{sec:ablation}

The ablation experiment attributes performance to the main components of
\method{}.  The full model combines the content coordinate, the structural
deviation coordinate, feature recovery, relation recovery, training
perturbation, and the set reader.  Each variant removes one part or replaces the
reader while preserving the same dataset row and metric as the main table.
These variants measure whether the point set gains its task value from
recoverable content, recoverable relations, relation conditioned deviation,
trainable adaptation, or set level aggregation.

Table~\ref{tab:ablation-results} reports the ablation results in the same
Emb/Result format as the main comparison, with the two quantities kept in
separate cells.  Removing feature recovery or relation recovery weakens many
downstream results, showing that approximate equivalence depends on preserving
both node content and pair relations.  Removing the structural deviation term
changes the point geometry most strongly on several graph tasks, while freezing
the point set reduces the ability to adapt the representation to the supervised
objective.  Reader substitutions mainly affect graph level rows, which is
consistent with the role of set aggregation in graph prediction.

\begin{table*}[t]
  \centering
  \caption{Ablation results from Experiment Result Ablation Table A format. Each variant keeps Emb and Result as separate columns.}
  \label{tab:ablation-results}
  \tiny
  \setlength{\tabcolsep}{1.4pt}
  \resizebox{\textwidth}{!}{%
  \begin{tabular}{llcccccccccccccccc}
    \toprule
    Dataset & Metric
    & \multicolumn{2}{c}{Full}
    & \multicolumn{2}{c}{\(-L_{\mathrm{feat}}\)}
    & \multicolumn{2}{c}{\(-\psi\)}
    & \multicolumn{2}{c}{\(-L_{\mathrm{link}}\)}
    & \multicolumn{2}{c}{\(-L_{\mathrm{emb}}\)}
    & \multicolumn{2}{c}{freeze-H}
    & \multicolumn{2}{c}{mean-pool}
    & \multicolumn{2}{c}{point-MLP} \\
    \cmidrule(lr){3-4}\cmidrule(lr){5-6}\cmidrule(lr){7-8}
    \cmidrule(lr){9-10}\cmidrule(lr){11-12}\cmidrule(lr){13-14}
    \cmidrule(lr){15-16}\cmidrule(lr){17-18}
    & & Emb & Result & Emb & Result & Emb & Result & Emb & Result
    & Emb & Result & Emb & Result & Emb & Result & Emb & Result \\
    \midrule
    HiGraph & F1 & \textbf{0.81} & \textbf{0.86} & 0.62 & 0.82 & 0.61 & 0.84 & 0.50 & 0.83 & 0.50 & 0.83 & 0.57 & 0.83 & 0.49 & 0.82 & 0.49 & 0.84 \\
    OCB & r & 0.84 & \textbf{0.90} & 0.46 & 0.62 & 0.51 & 0.63 & 0.53 & 0.64 & 0.54 & 0.64 & 0.35 & 0.50 & 0.89 & 0.62 & 0.87 & 0.62 \\
    Timing & R\textsuperscript{2} & \textbf{1.00} & \textbf{0.95} & 0.90 & 0.83 & 0.90 & 0.82 & 0.86 & 0.84 & 0.87 & 0.84 & 0.91 & 0.69 & -- & -- & 0.90 & 0.82 \\
    Photo & Acc & \textbf{0.88} & \textbf{0.95} & 0.86 & 0.92 & 0.85 & 0.93 & 0.85 & 0.93 & 0.85 & 0.93 & 0.81 & 0.72 & -- & -- & 0.86 & 0.92 \\
    Coauthor & Acc & \textbf{0.90} & \textbf{0.95} & 0.81 & 0.94 & 0.87 & 0.94 & 0.83 & 0.94 & 0.75 & 0.94 & 0.86 & 0.89 & -- & -- & 0.84 & 0.94 \\
    Photo int. & Acc & \textbf{0.97} & \textbf{0.94} & 0.86 & 0.92 & 0.85 & 0.93 & 0.85 & 0.93 & 0.85 & 0.93 & 0.81 & 0.72 & -- & -- & 0.86 & 0.92 \\
    Cora & Acc & \textbf{0.89} & \textbf{0.82} & 0.86 & 0.53 & 0.87 & 0.55 & 0.77 & 0.56 & 0.76 & 0.55 & 0.77 & 0.38 & -- & -- & 0.84 & 0.52 \\
    AIDS & Acc & \textbf{0.89} & \textbf{0.99} & 0.54 & 0.88 & 0.42 & 0.86 & 0.36 & 0.88 & 0.40 & 0.90 & 0.52 & 0.88 & 0.45 & 0.83 & 0.53 & 0.83 \\
    IMDB-M & Acc & 0.63 & \textbf{0.52} & 0.93 & 0.45 & 0.92 & 0.43 & 0.93 & 0.38 & 0.90 & 0.52 & 0.96 & 0.48 & 0.89 & 0.40 & 0.92 & 0.43 \\
    IMDB-B & Acc & \textbf{0.75} & \textbf{0.75} & 0.88 & 0.64 & 0.89 & 0.62 & 0.87 & 0.52 & 0.88 & 0.57 & 0.89 & 0.64 & 0.87 & 0.62 & 0.87 & 0.62 \\
    \bottomrule
  \end{tabular}}
\end{table*}

\subsection{WL Collision Experiment}
\label{sec:wl-diagnostic}

The WL collision experiment evaluates structural separation on controlled
nonisomorphic graph pairs.  The Weisfeiler Leman hierarchy provides the
standard reference for color refinement and higher order graph distinction
\cite{morris2019wl}.  BREC supplies graph pairs designed to expose realized
expressiveness in graph neural networks \cite{wang2023brec}.  For each pair,
\method{} trains the embedding and relation decoder with recovery objectives,
then compares whether the reconstructed graphs remain nonisomorphic after
consistent node permutations.

Table~\ref{tab:wl-results} reports the expanded BREC result from the workbook
using collision rate as the key metric.  On Basic and Extension, all covered
pairs pass direct graph recovery and the collision rate is zero.  The covered
set contains one hundred sixty official nonisomorphic pairs that are two
dimensional WL collisions, and their recovered graphs remain separated after
training.  A smaller sixty pair diagnostic gives the same separation behavior on
all recovery passing pairs: Basic ten out of ten, Extension ten out of ten, and
Regular three out of ten pass recovery, and the collision rate among the twenty
three passing pairs is zero.

\begin{table}[t]
  \centering
  \caption{WL collision rate from the BREC expanded 160-pair workbook sheet.}
  \label{tab:wl-results}
  \scriptsize
  \setlength{\tabcolsep}{4pt}
  \begin{tabular}{lcc}
    \toprule
    BREC category & Covered pairs & Collision rate \\
    \midrule
    Basic & 60 & \textbf{0.00} \\
    Extension & 100 & \textbf{0.00} \\
    \bottomrule
  \end{tabular}
\end{table}

These four experiments form a compact evidence chain.  The main experiment
shows downstream utility, the autoregressive experiment shows graph growth from
the same representation, the ablation experiment attributes the gains to the
recoverable set construction, and the WL experiment shows that the recovered
point sets separate the covered nonisomorphic collision pairs.
